%%%%%%%%%%%%%%%%%
% Configuration %
%%%%%%%%%%%%%%%%%

% Spec: 2-5 pages (excluding references). Do not change font, margins, or spacing.
\documentclass[12pt, letterpaper]{article}
\usepackage{xurl}
\usepackage[super,comma,sort&compress]{natbib}
\usepackage{amsmath}
\usepackage{graphicx}
\usepackage{listings}
\lstset{basicstyle=\ttfamily\small, columns=fullflexible}
\usepackage{lipsum}

% Any configuration that should be done before the end of the preamble:
\usepackage{hyperref}
\hypersetup{colorlinks=true, urlcolor=blue, linkcolor=blue, citecolor=blue}

\renewcommand{\refname}{References Cited}

\title{ECS 240 Project Proposal}
\author{Rajiv Choudhury, Jacob Hamzawi, Rishika Nigam, Gurpreet Singh}

\begin{document}

\maketitle

%%%%%%%%%%%
% Article %
%%%%%%%%%%%

\section{Introduction and Motivation}

\subsection{Problem Domain and Motivation}
% TODO: Application domain; what problem(s) we solve and why it's interesting.
% At least 3 citations, at least 2 of them research papers.

\lipsum[1]

\subsection{Proposed Language}
% TODO: Own toy language, or a subset of an existing one? What it contains/does,
% or which keywords/features we focus on. Include example syntax.

\subsection{Proposed Tool}
% TODO: Basic interpreter + one of: compiler, static analysis, dynamic analysis,
% type system. What are the inputs/outputs? How does a user interact with it?

\subsection{Scope and Related Work}
% TODO: Is this already solved/built? What existing tools solve similar problems?

\section{Goals}
% TODO: High level; details go in later sections.

\paragraph{Minimum goal.}
% Must have working by end of quarter.

\paragraph{Target goal.}
% Plan to have working by end of quarter.

\paragraph{Reach goal.}
% Extra features if time allows; nothing on the critical path.

\section{Syntax and Semantics}
% TODO: Example syntax, draft formal grammar, and at most a paragraph sketching
% the semantics (set of (INPUT, prog, OUTPUT) triples).

\begin{lstlisting}
example program goes here
\end{lstlisting}

\begin{align*}
  e &::= n \mid x \mid e_1 + e_2 \\
  c &::= x := e \mid c_1; c_2
\end{align*}

\section{Implementation Plan}

\subsection{Tool and Language Selection}
% TODO: What tools and languages will we use to build the project?

\subsection{Incorporating Existing Tools}
% TODO: From scratch, or using existing software/code bases? Justify each one.

\subsection{Architecture}
% TODO: Components and how they interact. Consider a diagram (can be tentative).

% \begin{figure}[h]
%   \centering
%   \includegraphics[width=\linewidth]{architecture.pdf}
%   \caption{Tool architecture.}
%   \label{fig:architecture}
% \end{figure}

\subsection{Development Effort}
% TODO: Functions/components/modules to build (refer to architecture).
% What is explicitly out of scope?

\subsection{Challenges}
% TODO: What could go wrong? What's the fallback if a step gets difficult?

\section{Collaboration Plan}
% TODO: How are components divided among members? Meeting plan (suggested
% 2x/week) and communication platform.

\paragraph{Timeline.}
% Expected weekly timeline of which components are developed when.

\begin{tabular}{ll}
  \textbf{Week} & \textbf{Components} \\
  \hline
  Week 1 & \\
  Week 2 & \\
\end{tabular}

\paragraph{Commitment statement.}

% TODO: uncomment when each person agrees
% \begin{itemize}
%   \item I commit to work and contribute equally to this project. --- \emph{Rajiv Choudhury}
%   \item I commit to work and contribute equally to this project. --- \emph{Jacob Hamzawi}
%   \item I commit to work and contribute equally to this project. --- \emph{Rishika Nigam}
%   \item I commit to work and contribute equally to this project. --- \emph{Gurpreet Singh}
% \end{itemize}

\section{Conclusion and Questions for Feedback}
% TODO: 1-2 paragraphs. What happens if successful, why it's exciting/important.
% What are we uncertain about / want feedback on?

%%%%%%%%%%%%%%
% References %
%%%%%%%%%%%%%%

\bibliographystyle{unsrtnat}
\bibliography{references}

% Material beyond the 5-page limit goes here (may or may not be graded):
% \appendix
% \section{...}

\end{document}
